
\documentclass{article}
\usepackage{colortbl}
\usepackage{makecell}
\usepackage{multirow}
\usepackage{supertabular}

\begin{document}

\newcounter{utterance}

\centering \large Interaction Transcript for game `cladder', experiment `full\_v1.5\_default', episode 2017 with Qwen3.6{-}27B.
\vspace{24pt}

{ \footnotesize  \setcounter{utterance}{1}
\setlength{\tabcolsep}{0pt}
\begin{supertabular}{c@{$\;$}|p{.15\linewidth}@{}p{.15\linewidth}p{.15\linewidth}p{.15\linewidth}p{.15\linewidth}p{.15\linewidth}}
    \# & \multicolumn{2}{c}{Player} && \multicolumn{2}{c}{Game Master} \\
    \hline

    \theutterance \stepcounter{utterance}  
    & & & \multicolumn{4}{p{0.6\linewidth}}{
        \cellcolor[rgb]{0.9,0.9,0.9}{
            \makecell[{{p{\linewidth}}}]{
                \texttt{\tiny{[P1$\langle$GM]}}
                \texttt{You are an expert in causal inference. The following question is not a typical commonsense query, but rather a meticulously designed question created by a professor specializing in causal inference, intended to assess the students' mastery of the course content.} \\
\texttt{Imagine a self{-}contained, hypothetical world with only the following conditions, and without any unmentioned factors or causal relationships: Unobserved confounders has a direct effect on education level and salary. Proximity to a college has a direct effect on education level. Education level has a direct effect on salary. Unobserved confounders is unobserved. For people living far from a college, the probability of high salary is 39\%. For people living close to a college, the probability of high salary is 49\%. For people living far from a college, the probability of college degree or higher is 57\%. For people living close to a college, the probability of college degree or higher is 30\%. Will college degree or higher increase the chance of high salary?} \\
\texttt{Start your answer with 'Yes' or 'No', followed by additional reasoning or evidence to support your explanation.} \\
            }
        }
    }
     \\ \\

    \theutterance \stepcounter{utterance}  
    & \multicolumn{4}{p{0.6\linewidth}}{
        \cellcolor[rgb]{0.561,0.737,0.561}{
            \makecell[{{p{\linewidth}}}]{
                \texttt{\tiny{[P1$\rangle$GM]}}
                \texttt{Yes} \\
\\ 
\texttt{**Reasoning:**} \\
\\ 
\texttt{To determine whether a college degree increases the chance of a high salary, we must estimate the **Causal Effect** of Education on Salary, not just the observed association. The problem explicitly states that there is an **unobserved confounder** that affects both Education and Salary. This creates a classic case of confounding bias.} \\
\\ 
\texttt{Let's define the variables:} \\
\texttt{*   \$Z\$: Proximity to college (Instrumental Variable candidate, though we must check assumptions).} \\
\texttt{*   \$X\$: Education level (Treatment).} \\
\texttt{*   \$Y\$: Salary (Outcome).} \\
\texttt{*   \$U\$: Unobserved confounder.} \\
\\ 
\texttt{The causal structure is:} \\
\texttt{1.  \$Z \textbackslash{}rightarrow X\$ (Proximity affects Education)} \\
\texttt{2.  \$X \textbackslash{}rightarrow Y\$ (Education affects Salary)} \\
\texttt{3.  \$U \textbackslash{}rightarrow X\$ and \$U \textbackslash{}rightarrow Y\$ (Unobserved confounder affects both)} \\
\\ 
\texttt{We are given marginal probabilities:} \\
\texttt{*   \$P(Y=1 | Z=0) = 0.39\$ (Far from college)} \\
\texttt{*   \$P(Y=1 | Z=1) = 0.49\$ (Close to college)} \\
\texttt{*   \$P(X=1 | Z=0) = 0.57\$ (Far from college)} \\
\texttt{*   \$P(X=1 | Z=1) = 0.30\$} \\
            }
        }
    }
    & & \\ \\

    \theutterance \stepcounter{utterance}  
    & & & \multicolumn{2}{p{0.3\linewidth}}{
        \cellcolor[rgb]{0.9,0.9,0.9}{
            \makecell[{{p{\linewidth}}}]{
                \texttt{\tiny{[GM$|$GM]}}
                \texttt{Yes} \\
            }
        }
    }
    & & \\ \\

    \theutterance \stepcounter{utterance}  
    & & & \multicolumn{2}{p{0.3\linewidth}}{
        \cellcolor[rgb]{0.9,0.9,0.9}{
            \makecell[{{p{\linewidth}}}]{
                \texttt{\tiny{[GM$|$GM]}}
                \texttt{no} \\
            }
        }
    }
    & & \\ \\

    \theutterance \stepcounter{utterance}  
    & & & \multicolumn{2}{p{0.3\linewidth}}{
        \cellcolor[rgb]{0.9,0.9,0.9}{
            \makecell[{{p{\linewidth}}}]{
                \texttt{\tiny{[GM$|$GM]}}
                \texttt{game\_result = LOSE} \\
            }
        }
    }
    & & \\ \\

\end{supertabular}
}

\end{document}
